% ======================================================================
%  KAVRAM: Reasoning  (Tier 1)  — kompakt
%  Chain of Thought / Reflection / Planning burada toplanır
% ======================================================================

\concept{Reasoning · Akıl Yürütme}{Reasoning}{Zor problemde modele cevaptan önce ``düşünme'' alanı vermek}

\stage{1}{Problem}
Çok adımlı, zor bir problemde modelden anında cevap istersen, yüzeysel ve sıklıkla
yanlış sonuç alırsın --- model ``patlatır''. Karmaşık mantık, aritmetik ve planlama,
ara adımlar olmadan tek hamlede doğru çözülmez.

\keyline{Tek hamlede cevap istemek, birine zor bir soruyu \emph{müsvedde kâğıdı vermeden} sormak gibidir.}

\stage{2}{Neden Var}
Modele cevaptan önce ara adımlar üretme alanı vermek --- ``yüksek sesle düşünmek'' ---
karmaşık görevlerde doğruluğu çarpıcı biçimde artırır. Reasoning desenleri bu
düşünmeyi yapılandırır: adım adım çözme, kendini eleştirip düzeltme, önce plan yapma.

\stage{3}{Tarihçe}
\begin{timeline}
\yr{2022} Chain of Thought: ``adım adım düşün'' \tlsep \yr{2022} ReAct: düşün → eyle → gözlemle \tlsep \yr{2023} reflection + planning: kendini eleştir, alt-adımlara böl \tlsep \yr{2024–26} reasoning modelleri: düşünme yerleşik
\end{timeline}

\stage{4}{Tanım}
\begin{problembox}
\textbf{Reasoning:} Modeli, nihai cevaba bağlanmadan önce ara adımlar üretmeye
yönlendiren teknikler. \textbf{Chain of Thought (CoT):} adım adım düşün.
\textbf{Reflection:} kendi çıktısını eleştir ve revize et. \textbf{Planning:}
eyleme geçmeden önce alt-adımlara böl.
\end{problembox}

\stage{5}{Zihinsel Model}
\begin{mentalbox}
Anında cevap istemek, birini \textbf{müsvedde olmadan} zor soruya yanıt vermeye
zorlamaktır. Reasoning ona şunları verir: bir \textbf{müsvedde kâğıdı} (CoT),
bir \textbf{ikinci bakış} (Reflection) ve başlamadan önce bir \textbf{plan}
(Planning). Yüksek sesle düşünmek modeli zor problemde akıllandırır --- ama
token ve gecikmeye mal olur. Kolay soruda müsvedde israftır.
\end{mentalbox}

\stage{6}{Mimari Diyagram}
Aynı soru, üç farklı yol. Sağa gittikçe daha çok ``düşünme'', daha çok maliyet.
\vspace{3pt}
\begin{center}
\begin{tikzpicture}[font=\sffamily\scriptsize, node distance=4mm,
   q/.style={draw=softgray, fill=white, rounded corners=6pt, minimum height=6.5mm,
             minimum width=11mm, align=center, inner sep=2pt},
   s/.style={draw=accent, fill=wash, rounded corners=2pt, minimum height=6.5mm,
             minimum width=13mm, align=center, inner sep=2pt},
   a/.style={draw=good, fill=good!10, rounded corners=2pt, minimum height=6.5mm,
             minimum width=11mm, align=center, inner sep=2pt},
   flow/.style={-{Stealth[length=1.6mm]}, accent, line width=0.7pt}]
  % Direct
  \node[q] (q1) {Soru};
  \node[a, right=8mm of q1] (a1) {Cevap};
  \draw[flow] (q1) -- (a1);
  \node[left=2mm of q1, font=\sffamily\bfseries\tiny, text=ink] {DIRECT};

  % CoT
  \node[q, below=7mm of q1] (q2) {Soru};
  \node[s, right=5mm of q2] (s2) {adım 1→2→3};
  \node[a, right=5mm of s2] (a2) {Cevap};
  \draw[flow] (q2) -- (s2); \draw[flow] (s2) -- (a2);
  \node[left=2mm of q2, font=\sffamily\bfseries\tiny, text=ink] {CoT};

  % Reflection
  \node[q, below=7mm of q2] (q3) {Soru};
  \node[s, right=5mm of q3] (d3) {taslak};
  \node[s, right=5mm of d3] (c3) {eleştir};
  \node[a, right=5mm of c3] (a3) {revize\\cevap};
  \draw[flow] (q3) -- (d3); \draw[flow] (d3) -- (c3); \draw[flow] (c3) -- (a3);
  \node[left=2mm of q3, font=\sffamily\bfseries\tiny, text=ink] {REFLECTION};

  \draw[-{Stealth[length=2mm]}, softgray, line width=0.7pt]
        ([xshift=6mm]a1.east) -- ++(0,-16mm)
        node[midway, right, font=\sffamily\tiny, text=softgray, align=left]
        {daha çok\\düşünme\\+ maliyet};
\end{tikzpicture}
\end{center}
\keyline{Reasoning bedavaya gelmez: her ara adım token ve gecikmedir. Doğruluğu kazandığın yerde harca.}

\stage{7}{Karar Çerçevesi}
Hangi desen, ne zaman? Görevin zorluğu ve hatanın maliyeti belirler.\par
\vspace{2pt}
\noindent\begin{tabularx}{\linewidth}{@{}>{\raggedright\arraybackslash}p{40mm} L@{}}
\arrayrulecolor{hairline}\toprule
{\sffamily\bfseries\small Görev} & {\sffamily\bfseries\small Desen} \\
\midrule
Basit arama / sınıflandırma & \textbf{Direct} (reasoning yok) \\
Çok adımlı mantık / aritmetik & \textbf{Chain of Thought} \\
Hata pahalı, kalite kritik & \textbf{Reflection} (eleştir-düzelt) \\
Araç dizisi / alt-hedef gerek & \textbf{Planning} \\
\bottomrule
\end{tabularx}

\stage{8}{Üretim Senaryosu}
\begin{prodbox}[Üretim Senaryosu · Yanlış Hesaplayan Asistan]
\textbf{Bağlam.} Bir finans asistanı çok adımlı hesaplarda (vergi, indirim, toplam)
sık sık yanlış sonuç veriyor --- doğru sayıları biliyor ama tek hamlede birleştirince
hata yapıyor.

\smallskip
\textbf{Çözüm.} CoT etkinleştirilir: model önce adımları yazar (ara toplamlar,
oranlar), sonra nihai sayıyı verir. Kritik akışlarda bir Reflection adımı
sonucu bir kez denetler.

\smallskip
\textbf{Sonuç.} Çok adımlı hata oranı belirgin düşer. Bedeli: biraz daha fazla
token --- yalnızca gerçekten zor akışlarda açıldığı için kabul edilebilir.
\end{prodbox}

\stage{9}{Anti-pattern}
\begin{antibox}[Anti-pattern · ``Her Yere Reasoning / Aşırı Düşünme'']
Basit görevlere bile CoT eklemek (boşa token ve gecikme); sonsuz reflection
döngüleri; ya da ham düşünme zincirini kullanıcıya göstermek. ``Daha çok reasoning
her zaman daha iyi'' sanmak.

\smallskip
\textbf{Belirti:} Basit cevaplar yavaş ve pahalı; bazen aşırı düşünüp basit işi bozuyor.\\
\textbf{Kök neden:} Reasoning'in bir maliyeti olduğunu görmemek.\\
\textbf{Düzeltme:} Reasoning'i zorluğa göre aç; basitte kapat; reflection'a sınır koy; ham zinciri kullanıcıdan gizle.
\end{antibox}

\stage{10}{Alternatifler \& Ödünleşimler}
\vspace{2pt}
\begin{tabularx}{\linewidth}{@{}>{\raggedright\arraybackslash}p{26mm} L L@{}}
\arrayrulecolor{hairline}\toprule
{\sffamily\bfseries\small Desen} & {\sffamily\bfseries\small En iyi} & {\sffamily\bfseries\small Maliyet} \\
\midrule
\textbf{Direct} & Basit, tek adımlı & En düşük \\
\addlinespace[2pt]
\textbf{Chain of Thought} & Çok adımlı mantık & Orta (ara adımlar) \\
\addlinespace[2pt]
\textbf{Reflection} & Yüksek kalite gereği & Yüksek (2× geçiş) \\
\addlinespace[2pt]
\textbf{Planning} & Karmaşık, araçlı görev & Yüksek; agent'a yaklaşır \\
\bottomrule
\end{tabularx}

\stage{11}{İlişkili Kavramlar}
\term{Agent} (planning + eylem döngüsü) \tlsep \term{Chain of Thought} \tlsep
\term{Reflection} \tlsep \term{Planning} \tlsep \term{Eval} (reasoning gerçekten iyileştirdi mi?)

\stage{12}{Özet}
\begin{itemize}
  \item Reasoning, modele cevaptan önce ``düşünme'' alanı verir; zor görevde doğruluğu artırır.
  \item CoT (adım adım), Reflection (eleştir-düzelt), Planning (önce böl).
  \item Bedavaya gelmez: her ara adım token ve gecikmedir.
  \item Zorluğa göre aç: basitte Direct, karmaşıkta CoT/Reflection/Planning.
  \item ``Daha çok reasoning her zaman daha iyi'' yanlıştır.
\end{itemize}

\stage{13}{Mühendislik İlkesi}
\begin{principlebox}
\textbf{İlke \#10 — Düşünmenin bir bedeli vardır; zor problemde öde, kolayda ödeme.}\\[4pt]
\normalsize Modele müsvedde kâğıdını ihtiyaç olduğunda ver. Gereksiz reasoning,
yokluğu kadar pahalıya patlar.
\end{principlebox}
